\part{Getting Started}

\chapter{Introduction}

FairWindSK is a touch-first marine display shell built around Signal~K.
It brings together live boat data, a suite of onboard web applications, and native
vessel-operation tools inside a single persistent frame designed for the helm.

\begin{notebox}
FairWindSK is part of the DYNAMO research programme at the University of Naples
'Parthenope', now supported by the DataX4Sea project funded by NEC Laboratories
America.
The goal is coastal environmental data crowdsourcing through a fleet of leisure
boats.
All results are open-source and open-data.
\end{notebox}

\section{What FairWindSK Does}

FairWindSK is a display \emph{shell}, not a data server.
It wraps a running Signal~K server, discovers the web applications installed on it,
and presents them in a stable, persistent frame alongside native operational controls.

Think of FairWindSK as the frame around your boat's apps and live navigation data.
It helps you keep key information visible while moving between tasks.

\subsection{Supported Platforms}

\rowcolors{2}{LtGray}{white}
\begin{longtable}{L{3.0cm}L{4.8cm}L{5.9cm}}
\toprule
\rowcolor{Navy}
\color{white}\sffamily\bfseries Platform &
\color{white}\sffamily\bfseries Typical use &
\color{white}\sffamily\bfseries Notes \\
\midrule
Raspberry Pi OS & Dedicated helm display, kiosk autostart &
  Full desktop feature set; kiosk helpers in \code{extras/} \\
Linux           & Navigation station laptop or PC    & Full desktop feature set \\
macOS           & Navigation station MacBook         & Full desktop feature set \\
Windows         & Navigation station PC              & Full desktop feature set \\
Android         & Helm tablet, secondary display     & Qt WebView mobile path \\
iOS / iPadOS    & iPad at helm or roving display     & Qt WebView mobile path \\
\bottomrule
\end{longtable}

\section{What FairWindSK Is NOT}

FairWindSK depends entirely on a running Signal~K server on your onboard network.
Without Signal~K it has no data to display, no apps to launch, and no instruments
to read.
It does not itself process NMEA sentences, connect directly to instruments, or store
charts.

\begin{warningbox}
FairWindSK is an aid to navigation.
It does not replace seamanship, lookout, paper charts, or certified safety
procedures.
The master of the vessel bears full responsibility for safe conduct of the passage
at all times.
\end{warningbox}

\section{What You Need Before Starting}

\begin{itemize}
  \item A running Signal~K server reachable on the boat's Wi-Fi network.
        The server must expose the standard API at
        \code{<server>/signalk} and the app catalogue at
        \code{<server>/signalk/v1/apps/list}.
  \item Your instruments (GPS, wind, depth, AIS, autopilot, engine monitor) feeding
        into Signal~K via NMEA~0183, NMEA~2000, or SeaTalk.
  \item FairWindSK installed on a display device.
\end{itemize}

\section{Device Guidance}

\subsection{Tablets and helm displays}

Tablets and larger marine displays are the best fit for primary helm use because
they provide larger touch targets, better layout stability, easier reading at helm
distance, and more room for simultaneous route and app awareness.

\subsection{Phones}

Phones work as companion displays.
Expect a denser layout and less room for simultaneous awareness.
For close pilotage, recovery operations, or high-workload situations, a larger
display is safer.

\subsection{Desktop systems}

On desktop platforms, FairWindSK can be used as a fixed navigation station
display.
The interface should be treated as a marine operational display first regardless
of input method.

\chapter{Before You Leave the Dock}

Run through these checks before every departure.
Five minutes at the dock is worth an hour at sea.

\section{First-Time Setup}

When FairWindSK is installed on a new display or freshly configured system, work
through these steps before relying on it underway:

\begin{enumerate}
  \item Open \ui{Settings~> Connection}.
        Enter the Signal~K server address, e.g.\ \code{http://192.168.1.100:3000}.
        Tap \key{Connect}.
        Watch the status change to \textit{Live}.
  \item Open \ui{Settings~> Units}.
        Choose knots for speed, nautical miles for distance, metres or feet for depth.
  \item Open \ui{Settings~> Applications}.
        Verify that the apps you expect are present.
        Arrange page~1 with the apps you use on every passage.
  \item Open \ui{Settings~> Comfort} and choose the \textit{Day} preset.
        Verify readability from your helm position.
  \item If your installation includes autopilot, anchor, alarms, or other
        plugin-backed features, confirm that each is available and responds correctly
        before departure.
  \item If the device has no physical keyboard, enable the virtual keyboard in
        \ui{Settings~> Main} and restart FairWindSK.
  \item Run a short test passage in calm conditions before relying on FairWindSK
        offshore.
\end{enumerate}

\begin{warningbox}
Do not assume a feature is ready simply because its icon is visible.
Confirm live data, control response, and real-world behaviour first.
\end{warningbox}

\section{Pre-Departure Checklist}

\rowcolors{2}{LtGray}{white}
\begin{longtable}{C{0.8cm}L{12.9cm}}
\toprule
\rowcolor{Navy}
\color{white}\sffamily\bfseries $\square$ &
\color{white}\sffamily\bfseries Check \\
\midrule
$\square$ & FairWindSK starts normally and the shell health state is normal \\
$\square$ & Signal~K connection is healthy (\textit{Live} in the status area) \\
$\square$ & Position live in Top Bar --- latitude and longitude plausible \\
$\square$ & SOG updating --- vessel has a valid GPS fix \\
$\square$ & Depth updating (if depth sounder is fitted) \\
$\square$ & Heading showing from compass or GPS \\
$\square$ & Chart plotter open and vessel in correct position on chart \\
$\square$ & Route or waypoint active if proceeding to a destination \\
$\square$ & AIS targets visible in chart plotter (if AIS receiver is fitted) \\
$\square$ & Display brightness and comfort preset match the current light conditions \\
$\square$ & Screen is readable and reachable from the normal operating position \\
$\square$ & POB button location confirmed --- know exactly where it is \\
$\square$ & Autopilot available and responding correctly (if required) \\
$\square$ & Anchor watch function confirmed available (if anchoring tonight) \\
\bottomrule
\end{longtable}
